\section{Introduction}

Container-based workloads have become a cornerstone of modern cloud infrastructure, revolutionizing how applications are developed, deployed, and scaled, particularly within microservices architectures and DevOps workflows. Unlike hypervisor-based virtual machines (VMs), which each require a full operating system and rely on a hypervisor for hardware virtualization, containers are lightweight, share the host system’s kernel, and incur minimal overhead. This efficiency makes them exceptionally well-suited for continuous integration and rapid scaling, allowing organizations to optimize resource utilization and reduce costs while ensuring quick, consistent deployments. Furthermore, containers enable seamless portability by packaging all userspace dependencies into a unified image, thereby avoiding the common problem of environment-specific discrepancies. Yet, despite containers’ numerous benefits, they also introduce new security challenges, especially where container confinement is concerned.

Linux kernel features such as cgroups, namespaces, and mount points provide resource, process, and network isolation for containers. However, this separation is not foolproof. Kernel vulnerabilities, such as buffer overflows or privilege escalation flaws can allow a process to circumvent the very mechanisms designed to confine it. Misconfigurations can similarly open holes in container security, for instance by granting unnecessary capabilities, mounting sensitive host directories, or running containers in privileged mode. Finally, sophisticated attacks like container escapes exploit these weaknesses to break out of the container’s resource and namespace isolation. Once an attacker gains elevated privileges or direct access to the host system, they can disrupt other containers, exfiltrate data, or compromise the entire infrastructure. 

Technologies like SELinux, AppArmor, and seccomp provide additional layers of security to mitigate these issues, but are notoriously complex to configure and prone to misconfiguration. A best practices guide for container security~\cite{CIS} can help with configuring a container, but it offers no guarantee that the given configuration will perform as intended by the end-user. As container adoption has skyrocketed across industries, particularly in multi-tenant cloud environments, these security challenges have become more pronounced.

To address the discussed challenges, the cloud industry has developed technologies like gVisor\footnote{\url{https://gvisor.dev}} and Kata Containers\footnote{\url{https://katacontainers.io}}, which isolate containers further by running them within lightweight VMs. While this approach is effective, gVisor-based containers bypass seccomp confinement by intercepting and emulating syscalls instead of directly blocking them through seccomp, which could result in security vulnerabilities if the emulation logic is flawed. By developing methods to audit and validate confinement mechanisms, container security could be simplified by reducing reliance on such solutions.

Test suites play a crucial role in verifying that systems meet specified requirements and function correctly under defined conditions. They are widely used across computing disciplines, from performance benchmarking to security assessment. For instance, the Standard Performance Evaluation Corporation (SPEC) \cite{henning2000spec} is best known for developing industry-standard benchmarks that evaluate CPU performance and power efficiency, while the Transaction Processing Performance Council (TPC) \cite{10.1007/978-3-642-18206-8_1} establishes benchmarking standards for database performance. The Phoronix Test Suite \cite{zboril2025performance} is a widely used benchmarking tool, particularly for Linux-based systems, providing automated performance evaluations across various hardware and software configurations. 

%Beyond performance testing, test suites are essential in software development to prevent regressions. Many organizations mandate that all code changes pass a comprehensive suite of tests before integration. In production environments, test suites help maintain reliability through uptime monitoring, ensuring that services remain operational. 
Specialized security-focused test suites play a vital role in detecting vulnerabilities in deployed systems. Tools such as network scanners aid in identifying misconfigurations, outdated software versions, and unintended service exposures, helping organizations proactively mitigate security risks. However, a significant gap remains in the context of containers. There are currently no standardized or widely adopted test suites designed specifically to evaluate container confinement, functionality, or isolation effectiveness. This absence leaves organizations to rely on ad hoc methods, which are often inadequate for uncovering subtle misconfigurations or ensuring robust confinement. Given the growing reliance on containerized environments in modern infrastructure, the development of comprehensive test suites tailored for containers is imperative. Such test suites would not only address these critical gaps but also provide a reliable foundation for improving trust and security in containerized systems.

In this paper, we describe Houdini, the first test suite for verifying container confinement. Given a Docker container configuration (including the kernel, docker, and runc version), Houdini will instantiate the container in a standalone QEMU-based VM and perform multiple tests (or \textit{tricks}, as we call them) to evaluate whether the container is susceptible to various forms of misconfigurations or failures. These tests explore vulnerabilities that could disrupt container functionality, such as
%, but not limited to 
improper resource allocation, inadequate isolation, or insecure filesystem setups. Houdini is written in Python and is easily extensible, providing an extension language for writing controlled test cases 
%(which we refer to as "tricks") 
designed to evaluate whether a container confinement mechanism properly enforces security boundaries.

%In this paper, we describe the motivation, design, and implementation of Houdini. We also present the results of case studies by showing how Houdini can be used to detect misconfigured containers. 

%maybe better for the conclusion
%Our hope with Houdini is that it will help with the deployment of better confined containers and will support the development of new technologies to more reliably confine containers.

The rest of this paper proceeds as follows. Section~\ref{sec:container-confinment-problem} defines container confinement and associated challenges. Section~\ref{sec:motivation} further motivates the need for a container security test suite. We describe Houdini's design and implementation in Sections~\ref{sec:design} and~\ref{sec:implementation}. Section~\ref{sec:evaluation} presents three case studies demonstrating how Houdini can detect security issues, misconfigurations, and other inconsistencies. We review related work in Section~\ref{sec:related} and conclude in Section~\ref{sec:conclusion}.